\chapter{Neck Lump}

\section*{Definition}
A neck lump is any palpable swelling or mass in the cervical region, with differential diagnosis guided by location, consistency, duration, and your age. Clinically significant lumps are those persistent beyond 2 weeks or larger than 1 cm.

\section*{What Happens in Your Body}
Neck masses arise from congenital remnants (branchial cleft, thyroglossal duct), reactive lymphoid hyperplasia from infection, neoplastic proliferation (primary or metastatic), or inflammation of salivary or thyroid tissue. The underlying mechanism involves cellular proliferation, fluid accumulation, or tissue hypertrophy within cervical anatomic compartments.

\section*{Causes}
\begin{itemize}
  \item \textbf{Congenital:} Branchial cleft cyst (anterior border of sternocleidomastoid), thyroglossal duct cyst (midline, elevates with tongue protrusion), cystic hygroma
  \item \textbf{Infectious:} Reactive lymphadenopathy (viral, bacterial, mycobacterial), suppurative lymphadenitis, abscess, cat-scratch disease, toxoplasmosis
  \item \textbf{Neoplastic:} Lymphoma, metastatic squamous cell carcinoma (especially from head/neck primaries), thyroid neoplasm, salivary gland tumor, lipoma, sebaceous cyst
  \item \textbf{Inflammatory:} Sj\"ogren's syndrome, sarcoidosis, Kawasaki disease, IgG4-related disease
\end{itemize}

\section*{When to See a Doctor}
See your doctor if:
\begin{itemize}
  \item Lump persists for more than 2--3 weeks without regression
  \item Rapid growth, firm or fixed consistency, or size greater than 2 cm
  \item Associated with unexplained weight loss, night sweats, or persistent fever
  \item Hoarseness, dysphagia, or otalgia (referred ear pain) are present
  \item History of smoking, heavy alcohol use, or prior head/neck cancer
\end{itemize}

\section*{Related Symptoms}
\begin{itemize}
  \item Sore throat, dysphagia, odynophagia, hoarseness, fever, weight loss, night sweats, otalgia
\end{itemize}
\textbf{Important:} In adults over 40, a persistent unilateral neck mass should be considered metastatic malignancy until proven otherwise, especially with a history of tobacco use.

\section*{Self-Care / Home Management}
\begin{itemize}
  \item If associated with upper respiratory infection, monitor for 2 weeks for spontaneous resolution
  \item Avoid unnecessary palpation or manipulation of the lump
  \item Warm compresses if tenderness suggests reactive lymphadenopathy
  \item Maintain good oral hygiene if source is suspected dental infection
\end{itemize}

\section*{How Your Doctor Will Evaluate You}
\begin{itemize}
  \item Detailed history (duration, growth pattern, associated symptoms) and physical exam (location, size, consistency, mobility, tenderness)
  \item Ultrasound with or without fine-needle aspiration (FNA) for cystic or solid characteristics
  \item Contrast-enhanced CT or MRI for deep-space masses or suspected malignancy
  \item Laboratory studies: CBC with differential, EBV/CMV serologies, thyroid function tests
  \item Excisional biopsy for persistent or suspicious solid masses after imaging
\end{itemize}

\section*{Treatment Options}
\begin{itemize}
  \item Observation and reassurance for self-limited reactive lymphadenopathy
  \item Antibiotics (amoxicillin-clavulanate or clindamycin) for bacterial lymphadenitis or abscess
  \item Surgical excision for thyroglossal duct cyst (Sistrunk procedure), branchial cleft cyst, or symptomatic benign tumors
  \item Referral to otolaryngology or oncology for confirmed or suspected malignancy
  \item Chemoradiation for metastatic head/neck squamous cell carcinoma; chemotherapy for lymphoma
\end{itemize}

\section*{References}
\begin{thebibliography}{99}
\bibitem{neckmass1} Pynnonen MA, Gillespie MB, Roman BR, et al. "Clinical practice guideline: evaluation of the neck mass in adults." Otolaryngol Head Neck Surg. 2017;157(2\_suppl):S1-S30.
\bibitem{neckmass2} Schwetschenau E, Kelley DJ. "The adult neck mass." Am Fam Physician. 2002;66(5):831-838.
\bibitem{neckmass3} Flint PW, Haughey BH, Lund VJ, et al. "Evaluation of neck masses." In: Cummings Otolaryngology. 7th ed. Elsevier; 2021.
\bibitem{neckmass4} Haynes J, Arnold KR, Aguirre-Oskins C, et al. "Evaluation of neck masses in adults." Am Fam Physician. 2015;91(10):698-706.
\bibitem{neckmass5} Bhattacharyya N. "Neck masses." In: Ballenger\'s Otorhinolaryngology. 18th ed. People\'s Medical Publishing House; 2016.
\bibitem{neckmass6} Celenk F, Baysal E, Durucu C, et al. "Differential diagnosis of neck masses." J Craniofac Surg. 2013;24(4):1399-1402.
\end{thebibliography}

\clearpage
